% !PRINTED 723
\weberpage{723}

Bezeichnen wir mit $\alpha$, $\beta$ die Indices von $t$, so ist $\chi_2(t) = (-1)^\alpha\Theta^\beta;$

\weq{2}{\overset{t}{\Sigma}\,\chi_2(t) = \Sigma\,(-1)^\alpha\Theta^\beta = \varphi(\Theta)}

ist eine \gesperrt{ganze Zahl des Körpers} $\Omega_\nu$, und es wird:
\[
X_2 = \frac{i\pi}{m}\,(-1,\Theta^{-1},r)\,\varphi(\Theta).
\]

Nun hat man für $\Theta$ alle Wurzeln der Gleichung
\weq{3}{\Theta^{\nu'} + 1 = 0}
zu setzen, und erhält für $A$, mit Rücksicht auf \S.\,197, (14):
\weq{4}{A = 2^{1-\nu'}\overset{\Theta}{\prod}\,\varphi(\Theta).}

Die Summe, durch die in (2) die Zahl $\varphi(\Theta)$ definirt ist,
enthält $\nu'$ Glieder von der Form $(-1)^\alpha\Theta^\beta$, worin $\alpha$, $\beta$ so zu
bestimmen sind, dass
\weq{5}{(-1)^\alpha\,5^\beta \equiv t \pmod{m}}
und $t$ zwischen 0 und $\nu$ liegt. Nun ist zunächst ersichtlich, dass
unter den Exponenten $\beta$ nicht zwei nach dem Modul $\nu'$ congruente vorkommen können. Denn ist $\beta \equiv \beta' \pmod{\nu'}$, so ist
$5^\beta \equiv 5^{\beta'} \pmod{\mu}$, und aus
\[
(-1)^\alpha\,5^\beta \equiv t, \quad (-1)^{\alpha'}\,5^{\beta'} \equiv t' \pmod{m}
\]
folgt:
\[
(-1)^\alpha\,t - (-1)^{\alpha'}\,t' \equiv 0 \pmod{\mu}.
\]

Da aber $t$ und $t'$ absolut kleiner als $\nu = \tfrac{1}{2}\mu$ sind, so ist
dies nur möglich, wenn $t = t'$, $\alpha = \alpha'$ ist. Demnach durchläuft
$\beta$ in (2) ein volles Restsystem nach dem Modul $\nu'$.

Da aber $\beta$ nach dem Modul $\nu$ zu nehmen ist, so werden die
nach (5) bestimmten $\beta$ theils kleiner, theils grösser als $\nu'$ ausfallen, und wenn wir also $\beta_1$ die Reihe der Zahlen $0, 1, \dots, \nu' - 1$
durchlaufen lassen, so erhalten wir in (2) zweierlei Arten von
Gliedern:
\weq{6}{\begin{aligned}
1.\ &(-1)^\alpha\Theta^\beta = (-1)^\alpha\Theta^{\beta_1}, & \beta = \beta_1 < \nu', \\
2.\ &(-1)^\nu\Theta^\beta = -(-1)^\alpha\Theta^{\beta_1}, & \beta = \beta_1 + \nu' \geqq \nu'.
\end{aligned}}
% [sic] in row 2 the page prints the exponent \nu where \alpha is required; reproduced as printed

Nun ist nach (5) der absolut kleinste Rest von $(-1)^\alpha\,5^\beta$
nach dem Modul $m$ positiv, und wenn wir daher eine Zahl
$c_\beta = \pm 1$ so bestimmen, dass der absolut kleinste Rest von $c_\beta\,5^\beta$
positiv wird, so ist in dem Falle (6), 1.
\[
c_{\beta_1} = (-1)^\alpha.
\]
